\begin{solution}
In harmonic gauge,
\[
\nabla^2\bar h_{00}=-16\pi G\rho.
\]
For a static Newtonian source the trace relation gives
\(h_{00}=-2\Phi\) with the sign convention
\(g_{00}=-(1+2\Phi)\), and the \(00\) Einstein tensor is
\(G_{00}=2\nabla^2\Phi\).  Thus
\[
G_{00}=8\pi G\rho
\quad\Longleftrightarrow\quad
\nabla^2\Phi=4\pi G\rho,
\]
fixing the normalization.
\end{solution}

\begin{solution}
For an isotropic fluid:
\[
\begin{array}{ll}
\text{null:}&\rho+p\geq0,\\
\text{weak:}&\rho\geq0,\ \rho+p\geq0,\\
\text{dominant:}&\rho\geq|p|,\\
\text{strong:}&\rho+p\geq0,\ \rho+3p\geq0
\quad(d=4).
\end{array}
\]
Dust and radiation satisfy all four; positive vacuum energy
\(p=-\rho\) saturates the null condition but violates the strong condition.
For \(\rho\geq0\) and \(p=w\rho\), the respective bounds are
\(w\geq-1\), \(-1\leq w\leq1\), and \(w\geq-1/3\) together with \(w\geq-1\).
\end{solution}

\begin{solution}
Tracing
\(R_{\mu\nu}-\tfrac12Rg_{\mu\nu}+\Lambda g_{\mu\nu}=0\) gives
\[
R=\frac{2d\Lambda}{d-2},\qquad
R_{\mu\nu}=\frac{2\Lambda}{d-2}g_{\mu\nu}.
\]
For a maximally symmetric space
\(R_{\mu\nu}=(d-1)K g_{\mu\nu}\), so
\[
K=\frac{2\Lambda}{(d-1)(d-2)}=\pm L^{-2}.
\]
Positive \(\Lambda\) gives de Sitter and negative \(\Lambda\) anti-de Sitter
curvature.
\end{solution}

\begin{solution}
For \(T_{00}=M\delta^3(\bm x)\),
\[
\nabla^2\bar h_{00}=-16\pi GM\delta^3(\bm x),
\qquad
\bar h_{00}=\frac{4GM}{r}.
\]
Trace reversal yields
\[
\dd s^2=-(1-2GM/r)\dd t^2
+(1+2GM/r)\dd\bm x^2.
\]
The time perturbation gives the equivalence-principle value \(2GM/b\) for
light deflection; the spatial perturbation supplies an equal amount.
\end{solution}

\begin{solution}
The linear field equation's divergence gives
\[
\partial^\mu h_{\mu\nu}-\alpha\,\partial_\nu h=0.
\]
Taking another divergence and the trace shows that for \(\alpha=1\),
\(h=0\) and then \(\partial^\mu h_{\mu\nu}=0\).  From ten components, four
transverse constraints and one trace constraint leave five massive spin-two
polarizations.  For generic \(\alpha\), the trace obeys its own second-order
equation and becomes a sixth scalar mode, with the wrong-sign kinetic term.
\end{solution}

\begin{solution}
Take a divergence and use
\(\nabla^\mu R_{\mu\nu}=\tfrac12\partial_\nu R\) and
\(\nabla^\mu T_{\mu\nu}=0\).  One obtains
\[
\partial_\nu(R+8\pi GT)=0,
\]
so \(R+8\pi GT=4\Lambda\) for a constant \(\Lambda\).  Substituting this trace
back into the trace-free equation gives
\[
G_{\mu\nu}+\Lambda g_{\mu\nu}=8\pi GT_{\mu\nu}.
\]
\end{solution}

\begin{solution}
The constraints are
\[
{}^{(3)}R+K^2-K_{ij}K^{ij}=16\pi G\rho+2\Lambda,\qquad
D_j(K^{ij}-h^{ij}K)=8\pi Gj^i.
\]
For the stated data only \({}^{(3)}R=0\) remains.  With
\(h_{ij}=\psi^4\delta_{ij}\),
\({}^{(3)}R=-8\psi^{-5}\nabla^2\psi\), so \(\nabla^2\psi=0\).
Asymptotic flatness and one puncture give
\[
\psi=1+\frac{M}{2r},
\]
the time-symmetric Schwarzschild slice in isotropic coordinates.
\end{solution}

\begin{solution}
The off-diagonal traceless spatial equation is
\[
\left(\partial_i\partial_j-\frac13\delta_{ij}\nabla^2\right)
(\Phi-\Psi)=8\pi G\,\Pi_{ij}.
\]
When anisotropic stress \(\Pi_{ij}\) vanishes and boundary modes are excluded,
\(\Phi=\Psi\).  Lensing probes \(\Phi+\Psi\), while slow-motion dynamics
mainly probes \(\Phi\); disagreement tests anisotropic stress or modified
gravity.
\end{solution}

\begin{solution}
For the Brans--Dicke form,
\[
\phi G_{\mu\nu}
=8\pi GT_{\mu\nu}
+\nabla_\mu\nabla_\nu\phi-g_{\mu\nu}\Box\phi
+\frac{\omega}{\phi}
\left(\nabla_\mu\phi\nabla_\nu\phi
-\frac12g_{\mu\nu}(\nabla\phi)^2\right).
\]
Combining its trace with the scalar variation gives
\[
(2\omega+3)\Box\phi=8\pi G\,T.
\]
Nonrelativistic matter has \(T\simeq-\rho\), whereas radiation has
\(T=0\), so only the former directly sources the scalar.
\end{solution}

\begin{solution}
Lovelock's theorem leaves
\[
\mathcal E_{\mu\nu}=aG_{\mu\nu}+b g_{\mu\nu}
\]
under the stated assumptions.  The apparent quadratic exception is
\[
\mathcal G=R_{\mu\nu\rho\sigma}R^{\mu\nu\rho\sigma}
-4R_{\mu\nu}R^{\mu\nu}+R^2.
\]
In four dimensions \(\int\sqrt{-g}\,\mathcal G\) is the Euler characteristic
plus a boundary term.  Its metric variation is therefore a total derivative
and contributes no bulk field equation.
\end{solution}
